\subsection{Модульная арифметика}
\textbf{Эйлер, Ферма, обобщённая Эйлера:} $a^{\varphi(m)}\equiv1\ (\gcd(a,m){=}1)$; для любых $a$ и $x\ge\log_2m$: $a^x\equiv a^{\varphi(m)+(x\bmod\varphi(m))}\pmod m$ (башни степеней). Вильсон: $(p{-}1)!\equiv-1\pmod p$; отсюда $n!\,(p{-}1{-}n)!\equiv(-1)^{n+1}\pmod p$ -- если $n$ близко к $p$ (например $p{-}1{-}n$ мало), $n!\bmod p$ считается через маленький факториал.
\par \textbf{Как пользоваться.} Деление по модулю -- умножение на обратный ($a^{-1}=a^{p-2}$ для простого $p$, либо расширенный Евклид для любого $m$ при $\gcd(a,m)=1$). Если $\gcd(a,m)>1$, делить нельзя: хранить числа как пару (степени простых, остаток) или работать с $\mathbb Z$ по КТО. Башня степеней $a^{b^c}\bmod m$ -- рекурсивно по $\varphi(m)$ (формула выше).
\par \textbf{Первообразные корни} есть при $m=1,2,4,p^k,2p^k$ ($p$ нечётное); их ровно $\varphi(\varphi(m))$. Порядок элемента делит $\varphi(m)$; $g$ -- корень $\iff g^{(p-1)/q}\ne1$ для всех простых $q\mid p{-}1$.
\par \textbf{Квадратичные вычеты:} $\left(\frac ap\right)\equiv a^{(p-1)/2}$; взаимность $\left(\frac pq\right)\left(\frac qp\right)=(-1)^{\frac{p-1}2\cdot\frac{q-1}2}$ (нечётные простые $p\ne q$). Для нечётного $p$ и $a\not\equiv0$ корней $x^2{=}a$ ровно $0$ или $2$.
\par \textbf{$\varphi$:} $\sum_{d\mid n}\varphi(d)=n$; $\varphi(mn)=\varphi(m)\varphi(n)\frac{g}{\varphi(g)}$, $g=\gcd(m,n)$; $\sum_{k\le n,\gcd(k,n)=1}k=\frac{n\varphi(n)}2$ ($n>1$).
\par \textbf{Диофант} $ax+by=c$: решения есть $\iff\gcd(a,b)\mid c$; все: $(x_0+k\frac b g,\ y_0-k\frac ag)$. При $\gcd(a,b)=1$ число решений в $x,y\ge0$: $\frac n{ab}-\{\frac{b'n}a\}-\{\frac{a'n}b\}+1$, где $aa'\equiv1\ (\mathrm{mod}\ b)$, $bb'\equiv1\ (\mathrm{mod}\ a)$.
\par \textbf{Степени и НОД:} $\gcd(a^n-1,a^m-1)=a^{\gcd(n,m)}-1$; $v_p(n!)=\sum_{i\ge1}\lfloor n/p^i\rfloor=\frac{n-s_p(n)}{p-1}$ ($s_p$ -- сумма $p$-ичных цифр; число нулей в конце $n!$ -- это $v_5$).
\par \textbf{НОД/НОК:} $\gcd(a,\mathrm{lcm}(b,c))=\mathrm{lcm}(\gcd(a,b),\gcd(a,c))$ и двойственное; $\gcd(\mathrm{lcm}(a,b),\mathrm{lcm}(b,c),\mathrm{lcm}(a,c))=\mathrm{lcm}(\gcd(a,b),\gcd(b,c),\gcd(a,c))$.
\subsection{Мёбиус и суммы по делителям}
$\mu(n)=0$ при квадрате в $n$, иначе $(-1)^{\#\text{простых}}$; $\sum_{d\mid n}\mu(d)=[n=1]$.\\
$g(n)=\sum_{d\mid n}f(d)\iff f(n)=\sum_{d\mid n}\mu(d)g(\tfrac nd)$;\ \ $g(n)=\sum_{n\mid d}f(d)\iff f(n)=\sum_{n\mid d}\mu(\tfrac dn)g(d)$;\\
$g(n)=\sum_{m\le n}f(\lfloor\tfrac nm\rfloor)\iff f(n)=\sum_{m\le n}\mu(m)g(\lfloor\tfrac nm\rfloor)$.
\par \textbf{Число делителей и их сумма:} $d(n)=\prod(e_i{+}1)$, $\sigma(n)=\prod\frac{p_i^{e_i+1}-1}{p_i-1}$ по разложению $n=\prod p_i^{e_i}$; суммы по всем $k\le n$ считаются перестановкой суммирования: $\sum_{k\le n}d(k)=\sum_{k\le n}\lfloor\tfrac nk\rfloor$, $\ \sum_{k\le n}\sigma(k)=\sum_{k\le n}k\lfloor\tfrac nk\rfloor$.
\par \textbf{Взаимно простые пары.} $\#\{(i,j)\le n:\gcd(i,j)=1\}=\sum_{d\le n}\mu(d)\lfloor\tfrac nd\rfloor^2=2\sum_{i\le n}\varphi(i)-1$; пар с $\gcd=k$ -- то же для $\lfloor n/k\rfloor$. Это основной приём для «посчитать пары с условием на НОД»: перейти к $\mu$ и сгруппировать по $d$.
\par \textbf{Блоки $\lfloor n/i\rfloor$.} Значение $\lfloor n/l\rfloor$ постоянно на $[l,\ r]$, $r=\lfloor n/\lfloor n/l\rfloor\rfloor$: цикл \texttt{for(l=1;l<=n;l=r+1)\{r=n/(n/l);...\}} делает $O(\sqrt n)$ шагов. Применять к любым суммам вида $\sum_i f(\lfloor n/i\rfloor)\,g(i)$ с префиксными суммами $g$.
\par Свёртки: $\sum_{d\mid n}\varphi(d)=n$, $\ \sum_{k\le n}\gcd(k,n)=\sum_{d\mid n}d\varphi(\tfrac nd)$, $\ \sum_{k\le n}\mathrm{lcm}(k,n)=\frac n2\big(\sum_{d\mid n}d\varphi(d)+1\big)$,\\
$\sum_{i,j\le n}\mathrm{lcm}(i,j)=\sum_{l\le n}S\big(\lfloor\tfrac nl\rfloor\big)^2\sum_{d\mid l}\mu(d)\,l\,d$, где $S(m)=\frac{m(m+1)}2$.\\
Бесквадратных $\le N$: $\sum_{k\le\sqrt N}\mu(k)\lfloor N/k^2\rfloor$. Значений $\lfloor n/i\rfloor$ не более $2\sqrt n$ (блоки одинаковых значений); $\sum_{k\le n}\lfloor n/k\rfloor\approx n\ln n$; $\sum_{d\mid n}d=O(n\log\log n)$.
\par \textbf{Идея: суммы по делителям и свёртки.} Хотим $g(m)=\sum_{d\mid m}f(d)$ для всех $m\le n$ -- «префиксные суммы» по решётке делимости. Делаем по одному простому $p$: для $i=1..n/p$ по возрастанию $f[ip]\mathrel{+}=f[i]$ (после всех простых получается $g$), всего $O(n\log\log n)$. Обратное -- \emph{преобразование Мёбиуса}: те же простые, но $i$ по убыванию и $f[ip]\mathrel{-}=f[i]$ (это и есть $f=\mu*g$). Для сумм по кратным ($g(d)=\sum_{d\mid m}f(m)$) циклы идут в обратную сторону. Свёртки: \textbf{НОД}: $c_k=\sum_{\gcd(i,j)=k}a_ib_j$ -- сумма по кратным $\to$ почленное произведение $\to$ Мёбиус по кратным. \textbf{НОК}: то же по делителям. Свёртка Дирихле $\sum_{ij=n}a_ib_j$ -- за $O(n\log n)$ прямым перебором; $\mu*\mathbf1=[n{=}1]$, $\varphi*\mathbf1=\mathrm{id}$. Код: \texttt{math/dirichlet.cpp}.
\subsection{Суммы квадратов, тройки, решётка}
\textbf{Тройки:} $a=k(m^2{-}n^2),\ b=2kmn,\ c=k(m^2{+}n^2)$, $m>n>0$, $\gcd(m,n)=1$, $m\not\equiv n\ (2)$ -- все примитивные (при $k{=}1$). Пар $(a,b)$, $a,b>0$, с $a^2{+}b^2=n^2$: $\prod_{p\equiv1\,(4)}(2\alpha_p{+}1)-1$.
\par $n$ -- сумма двух квадратов $\iff$ простые $\equiv3\ (4)$ входят в чётных степенях; $r_2(n)=4(d_1(n){-}d_3(n))$ (делители $\equiv1,3$ mod 4); любое $n$ -- сумма четырёх квадратов, $r_4(n)=8\sum_{d\mid n,\,4\nmid d}d$; $r_8(n)=16\sum_{d\mid n}(-1)^{n+d}d^3$; сумма трёх квадратов невозможна ровно для $n=4^a(8b{+}7)$.
\par \textbf{Пик:} $S=I+\frac B2-1$; целых точек на отрезке $\gcd(|dx|,|dy|){+}1$.
\subsection{Фибоначчи}
$F_{2n}=F_n(2F_{n+1}{-}F_n)$,\ $F_{2n+1}=F_n^2+F_{n+1}^2$;\ $\gcd(F_a,F_b)=F_{\gcd(a,b)}$;\ Кассини $F_{n-1}F_{n+1}-F_n^2=(-1)^n$;\ $\sum_{i\le n}F_i=F_{n+2}-1$,\ $\sum F_i^2=F_nF_{n+1}$. $n$ -- число Фибоначчи $\iff5n^2\pm4$ -- точный квадрат. Цекендорф: единственное разложение без соседних слагаемых. Период Пизано $\pi(m)\le6m$, равенство при $m=2\cdot5^k$, $k\ge1$.
\subsection{Асимптотики}
$n!\sim\sqrt{2\pi n}\big(\frac ne\big)^n\big(1+\frac1{12n}+\frac1{288n^2}-\frac{139}{51840n^3}\big)$;\ $\ln n!=n\ln n-n+\frac12\ln(2\pi n)+\frac1{12n}-\frac1{360n^3}+\frac1{1260n^5}$;\ $\binom{2n}n\approx\frac{4^n}{\sqrt{\pi n}}$;\ $H_n\approx\ln n+0.5772157+\frac1{2n}$;\ $\pi(x)\approx x/\ln x$.
\subsection{Таблицы}
\begin{center}\footnotesize\setlength{\tabcolsep}{2.5pt}
\begin{tabular}{c|cccccccccccc}
$k$&1&2&3&4&5&6&7&8&9&10\\\hline
$\pi(10^k)$&4&25&168&1229&9592&78498&664579&5761455&50847534&455052511\\
\end{tabular}\end{center}
\par Наибольшее число делителей ($N<$): $5\cdot10^4$ -- 100 ($45360$); $10^7$ -- 448 ($8648640$); $10^{10}$ -- 2304 ($6983776800$)\,$^\dagger$; $10^{19}$ -- 161280.
\begin{center}\footnotesize\setlength{\tabcolsep}{2.5pt}
\begin{tabular}{c|cccccc}
$N\le$&$10^3$&$10^6$&$10^9$&$10^{12}$&$10^{15}$&$10^{18}$\\\hline
$n$&840&720720&735134400&963761198400&866421317361600&897612484786617600\\
$d(n)$&32&240&1344&6720&26880&103680\\
\end{tabular}\end{center}
\par \textbf{Пределы типов:} $12!<2^{31}\le13!$; $20!<2^{63}\le21!$; $170!<$ DBL\_MAX $<171!$; $F_{46}<2^{31}$, $F_{92}<2^{63}$; $C_{19}<2^{31}$, $C_{35}<2^{63}$ (Каталан); $\binom{66}{33}<2^{63}$.
\par \textbf{Простые.} NTT: $998244353=119\cdot2^{23}{+}1$, $167772161=5\cdot2^{25}{+}1$, $469762049=7\cdot2^{26}{+}1$ (корень 3), $962592769=459\cdot2^{21}{+}1$; 64-бит: $29\cdot2^{57}{+}1$ ($g{=}3$), $27\cdot2^{56}{+}1$ ($g{=}5$), $237\cdot2^{52}{+}1$ ($g{=}5$).
$10^9{+}\{7,9,21,33,87,93,97,103,123\}$;\ $10^{18}{+}\{3,9,31,79,177,183,201,283\}$. Для хеша: $970592641$ (30 бит)$^\dagger$, $31443539979727$ (45), $3006703054056749$ (52), $1791791791$, $1234676767$, $1000575757$, $998696969$, $996179179$, $991676767$.
\par \textbf{Миллер--Рабин, детерминированные базы:} $n<4759123141$: $\{2,7,61\}$; $n<2152302898747$: $\{2,3,5,7,11\}$; $n<2^{64}$: $\{2,325,9375,28178,450775,9780504,1795265022\}$ или первые 12 простых.
